% !TEX root = ../main.tex
\chapter{논의}
\label{ch:discussion}

\ref{ch:results}장에서는 아비트럼 원(Arbitrum One)에서 잰 결과를 보고했어요. 이 장에서는 \ref{ch:introduction}장의 연구 질문(RQ)과 목표(O)를 하나하나, 그 장에서 밝힌 기준에 맞추어 따져 봐요. 그리고 모든 해석마다 그 해석이 기대는 \ref{ch:results}장의 표를 가리켜요. 따로 말하지 않으면, 같은 기간(점수를 계산한 바로 그 기간) 안에서 한 이야기는 매칭 코호트(모든 점수가 똑같이 점수를 매기는 지갑 묶음)를 두고 한 것이에요. 그리고 홀드아웃(나중 기간을 떼어 두고 맞혀 보는 시험)에서 한 이야기는 스펜더(spender, 허락을 받아 대신 쓰는 쪽) 코호트를 두고 한 것이에요.

\dissertationSubheading[sec:findings-rq]{연구 질문별로 알아낸 것}

\dissertationSubsubheading[sub:finding-rq1]{RQ1: 평판 구조와 순위 갈라짐}

RQ1은 엔도스랭크(EndorseRank)와 AWP(적응형 가중 페이지랭크)가 같은 지갑들을 서로 다른 순서로 줄 세우는지 물어요. 답은 “그렇다”예요. 두 점수가 얼마나 같은 순서인지 나타내는 값의 신뢰구간(믿을 만한 범위)에는 0도 1도 들어 있지 않아요(표~\ref{tab:tie-sensitivity}). 두 점수 사이에 있는 일치는 그래프(점과 화살표로 이은 그림)보다 처리 과정(파이프라인)의 규칙 하나에서 더 많이 나와요. 엔도스랭크는 지갑 쌍 대부분을 동점으로 만들어요. 그리고 엔도스랭크가 순서를 정하는 쌍은 대부분, 허락(allowance, 토큰을 얼마까지 대신 써도 된다는 허락) 그래프 안의 지갑 하나를 그래프 밖의 지갑 하나와 맞세운 쌍이에요. 그래프 밖의 지갑에 0점을 주는 대신 외톨이 점(화살표가 하나도 없는 점)으로 남겨 두면, 그 지갑들은 공통 값의 동점 덩어리에 들어가요. 그러면 일치는 0보다 조금 아래로 떨어져요(\ref{sec:rank-divergence}절).

이 순위 갈라짐은 두 그래프를 만드는 방법 때문에 생겨요. 보증(믿고 맡김) 그래프는 마지막 허락(가장 최근 허락)을 기록해요. 송금(transfer) 그래프는 모든 송금을 기록하는데, 보증 그래프의 화살표(엣지) 수는 송금 그래프의 약 9분의 1이에요\cend{3}. 그리고 두 방법은 다시 출발하는(재시작) 방식도 달라요. 페이지랭크(PageRank)는 주어진 화살표에 따라 정상 분포의 몫(오래 걸었을 때 각 점에 머무는 비율)을 나누어 놓아요\cend{15,23}. 그래서 활발한 송신자(보낸 쪽)에게서 송금은 많이 받지만 허락은 적게 받은 지갑은, AWP에서는 순위가 높고 엔도스랭크에서는 순위가 그리 높지 않아요. 많은 주인(토큰 주인)이 허락해 주지만 송금 허브(송금이 많이 모이는 곳)는 아닌 계약(스마트 계약, 블록체인 위에서 저절로 실행되는 프로그램)은 그 반대예요. 가설 H1은 뒷받침돼요. 실제로는 한 점수가 다른 점수를 대신할 수 없어요. 빠르다는 이유로 엔도스랭크를 고른 운영자는 AWP를 더 빨리 어림한 값이 아니라, 아예 다른 순위를 얻게 돼요.

\dissertationSubsubheading[sub:finding-rq2]{RQ2: 같은 기간 일치도}

RQ2는 각 점수가 같은 시기의 허락 점검, 송금 점검과 얼마나 잘 맞는지 물어요. 점수는 저마다 자기가 걷는 화살표를 따라가요. 송금 그래프를 걷는 점수들은 송금 지표 가족(비슷한 지표 묶음)과 맞고, 허락 지표 가족에서는 0 가까이에 머물러요. 허락 그래프를 걷는 점수들은 허락 지표 가족과 맞아요. 그리고 각 점수가 자기 지표 가족에서 얻은 구간은 모두 0을 벗어나요(표~\ref{tab:alignment-hybrid}).

H2에는 두 가지 단서가 붙어요. 첫째는 순환성(같은 자료로 만든 것끼리 견주는 문제)이에요. 한 지표 가족에 든 대리 지표(대신 재는 값)는 모두, 짝이 되는 그래프와 같은 이벤트(계약이 남기는 기록)로 계산해요. 그래서 이 일치는 Cronbach와 Meehl\cend{22}이 말한 뜻의 구성 타당도(점수가 재려는 것을 정말 재는지)이지, 외부 검증(바깥 자료로도 맞는지)이 아니에요. 둘째는 엔도스랭크의 동점 구조예요(\ref{sec:rank-divergence}절). 엔도스랭크의 허락 계수(두 순위가 얼마나 같은지 나타내는 수)는 허락을 받는 지갑 131개에 기대고, 동점 때문에 오를 수 있는 천장(얻을 수 있는 가장 큰 값)에 닿아 있어요. 송금 계수는 대부분, 어떤 지갑이 허락 그래프에 나타나기라도 하는지를 보여 줄 뿐이에요. 두 계수는 천장이 서로 달라요. 그래서 이 두 계수로는 엔도스랭크가 송금보다 허락에 더 기대는지에 대해 아무것도 말할 수 없어요. 외톨이 점 규칙을 쓰면 두 계수가 모두 크게 바뀌어요(표~\ref{tab:tie-sensitivity}). 하지만 이 규칙은 결과를 안 뒤에 시험해 본 것이라서, H2는 이 규칙에 기대지 않아요. 점수가 자기 화살표로 만든 점검과 맞는다는 좁은 뜻에서는 H2가 뒷받침돼요. 이 비교들은 어느 것도 미리 정해 두지 않았어요. 엔도스랭크와 AWP 사이의 짝지은 차이(같은 지갑들 위에서 두 점수를 짝지어 견준 차이)는 부록~\ref{ch:appendix-er-awp}에 있고, 이 차이를 근거로 한 주장(C)은 없어요.

\dissertationSubsubheading[sub:finding-rq3]{RQ3: 계산 비용}

RQ3은 같은 표본에서 각 점수를 계산하는 데 얼마나 드는지 물어요. 비용은 화살표 수를 따라가요. 엔도스랭크는 시간, 메모리, 반복(한 바퀴) 횟수 모두에서 가장 싼 점수예요. 송금 층(송금 화살표로 만든 그래프)을 걷는 두 혼합 점수는 AWP보다 오래 걸려요(표~\ref{tab:benchmark-runtime}). 시간을 잰 호출에서는 행렬을 짜 맞추는 일이 시간을 가장 많이 차지해요. 짜 맞추기와 반복은 둘 다 $|E|$에 비례해서 늘어나요\cend{23}. 그래서 시간 비율은 화살표 수의 비율을 따라가고, 반복 횟수의 차이는 별로 보태는 것이 없어요. 엔도스랭크의 그래프가 더 작은 까닭은, AWP는 모든 송금을 담지만 엔도스랭크는 마지막 허락들을 한순간 찍은 모습(스냅숏)만 담기 때문이에요\cend{3}. 두 점수는 똑같은 계산을 돌려요.

규모에 따른 변화를 보인 표들은 이 결과를 넓혀 주기보다, 이 결과에 단서를 달아요(표~\ref{tab:benchmark-scaling}). 코호트 밖의 확장 지갑 풀(더 큰 주소 모음) 주소들에는 꺼낸 화살표가 하나도 없어요. 그래서 어느 단계에서도 코호트 그래프보다 큰 그래프는 시험하지 못해요. 감쇠 계수(화살표를 따라갈 확률)를 바꾼 점검과 상위 토큰만 남긴 점검은 일치도의 모양과 실행 시간 비율을 바꾸지 않아요(표~\ref{tab:robustness-damping}과~\ref{tab:robustness-tokens}). H3은 매칭 코호트와 그 부분 그래프들에서 뒷받침돼요.

\dissertationSubsubheading[sub:finding-rq4]{RQ4: 앞으로 생길 새 허락에 대한 시간 홀드아웃}

RQ4에는 두 부분이 있어요. 첫째 부분은 2월 말에 고정한 점수가, 그다음 석 달 동안 스펜더 코호트에서 생긴 새 주인--스펜더 허락과 새 송금자(새로 송금을 보내온 주소)를 예측하는지(미리 맞히는지) 물어요. 모든 라벨(나중에 채점에 쓰는 정답 값)은 점수에 쓴 이벤트보다 뒤에 생긴 것이에요. 그래서 이것은 같은 기간 안의 점검이 아니에요. 대부분의 점수는 정말로 두 라벨을 모두 예측해요. 그리고 얼마나 잘 맞히는지는 화살표의 종류가 정해요(표~\ref{tab:holdout-spenders}). 새 허락 쌍(새로 생긴 주인--스펜더 허락)에서는 받은 허락 수(허락 입차수)가 앞서요. 그 뒤를 허락 그래프를 걷는 점수들과 중간 비율로 섞은 혼합 점수들이 따르고, 송금만으로 만든 점수들은 낮게 머물러요. 새 송금자에서는 송금 입차수(송금을 보내온 서로 다른 주소 수)가 앞서요. 하지만 많은 주인이 허락해 준 스펜더는 새 송금자도 많이 얻어요. 그래서 받은 허락 수와 허락 걷기가 송금 그래프를 걷는 어떤 점수보다도 이 라벨을 더 잘 줄 세워요. 이 증거는 스펜더에 대한 것이에요. 매칭 코호트의 트레이더(거래하는 사람)에서는 송금으로 만든 점수들만 나중에 이익을 내며 포지션(열어 둔 거래)을 닫은 횟수를 줄 세웠어요. 그리고 이익을 내며 닫은 비율은 어떤 점수도 줄 세우지 못했어요(\ref{sub:trader-results}절).

둘째 부분은 어떤 페이지랭크가 자기가 매끄럽게 다듬는(평활하는) 원시 차수(화살표를 그냥 센 수)를 이기는지 물어요. 어느 기간에서도 이긴 페이지랭크는 없어요. 짝지은 차이는 모두 0보다 한참 아래에 있어요(표~\ref{tab:holdout-diff}과~\ref{tab:fresh-posthoc}). 그리고 AWP가 차수에 못 미치는 것은, 동결일(점수를 고정하는 날) 전에 송금을 받은 적이 있는 스펜더들에서도 그대로예요. 그러니 AWP가 0점을 준 스펜더들로는 이것을 설명할 수 없어요(표~\ref{tab:holdout-receiving}). 한 지갑이 한 분기(석 달) 안에 새로 얻는 거래 상대의 수는, 그 지갑이 이미 가진 거래 상대의 수와 상관이 있어요(함께 움직여요). 그리고 지갑 바로 곁의 화살표만 센 차수는 이 정보를, 퍼뜨림(전파)이 끌어들이는 묽어짐 없이 담아요. 엔도스랭크는 각 주인의 보증을 그 주인이 허락해 준 모든 스펜더에게 나누어 주고, AWP는 각 송신자의 것을 그 송신자가 돈을 보낸 모든 주소에게 나누어 주기 때문이에요. 그래서 페이지랭크가 기간 밖에서 잡아내는 것은 무엇이든 자기가 다듬는 차수에서 물려받은 것이고, 차수에 더 보태지 않아요. 그렇다고 첫째 부분이 뒤집히지는 않아요. 그 개수들도 같은 화살표의 통계이기 때문이에요. 이 결과가 배제하는 것은, 이 라벨들에서 퍼뜨림이 차수를 넘어서는 예측력을 보탠다는 가능성이에요. C-PR ($\lambda=1$)의 증가분(자기 원시 차수보다 나은 정도)을 빼면, 이 대비(두 점수를 견준 차이)들은 라벨을 안 뒤에 계산한 것이에요.

RQ4는 연구 계획서와 달라요. 계획서는 점수를 거래 성공과 견주어 검증하고, 그 검증의 한계를 결과에 바탕을 둔 기준 점수로 정했어요. 그 기준 점수는 순환적(같은 자료를 돌려 쓴 것)이었고, 그래서 거두어들였어요(\ref{sub:gf-pr}절). 대신 쓴 라벨들은 점수 기간 밖에 있지만 같은 구성개념(점수가 재려고 하는 생각)에 속해요. 그리고 두 라벨 어디에도 이익이나 빚 갚기의 요소는 전혀 없어요. 그래서 RQ4는 각 구성개념의 시간 예측 타당도(시간이 지난 뒤의 일을 얼마나 잘 미리 맞히는지)를 보여 주고, 이때 원시 차수가 더 강한 기준선(비교 기준)이에요. 신용 위험은 다루지 않고 남겨 둬요. 등록해 둔(계획을 먼저 공개해 둔) 기간의 청산(손실이 너무 커져 거래가 강제로 끝나는 일) 라벨이 신용 위험에 가장 가까워요(\ref{sub:fresh-holdout}절).

\dissertationSubsubheading[sub:finding-rq5]{RQ5: 두 종류의 화살표 합치기}

RQ5의 두 부분은 모두, 저마다의 규칙 아래에서 “아니요”라는 답을 얻어요. 둘째 부분은 두 층을 한 번에 걷는 걷기가, 한 층만 걷는 것 어느 쪽보다 지갑을 더 잘 줄 세우는지 물어요. 9월 14일에 정한 규칙으로는 그렇지 않아요. 봄 홀드아웃에서 C-PR(결합 페이지랭크, 두 그래프를 함께 걷는 점수)은 새 허락 쌍에서 두 걷기 사이에 놓이고, 새 송금자에서 송금 걷기를 이긴다는 것이 보이지 않아요(표~\ref{tab:holdout-diff}). 그리고 같은 기간 안에서는 송금 걷기를 따라가고, 허락 지표 가족에서는 0 가까이에 머물러요. 그래서 보조 규칙은 허락 부분에서 실패해요(\ref{sec:coupled-results}절). 첫째 부분은 이 걷기가 송금만 걷는 걷기보다 나은지 물어요. 봄 홀드아웃에서 한 탐색적(미리 정하지 않고 살펴본) 비교는, 새 허락 쌍에서는 낫고 새 송금자에서는 알아챌 만한 손실이 없다고 했어요. 하지만 이 비교는 라벨을 안 뒤에 한 것이에요. 이 질문은 등록 재현(계획을 먼저 공개해 두고 새 자료로 다시 해 보기)이 결정하고, 그 답은 “아니요”예요. C-PR은 이번에도 새 허락 쌍에서 송금 걷기를 이겼고, 청산 없이 포지션을 닫은 트레이더를 조금 더 잘 줄 세웠어요. 하지만 새 송금자에서 비열등성(정한 폭보다 더 나쁘지 않음)은 보이지 못했고, 규칙은 세 조건을 모두 채워야 했어요(\ref{sub:fresh-results}절). 등록한 민감도 분석(조건을 바꿔 결과가 얼마나 움직이는지 보기)의 비교 대상인 AWP와 견주면, C-PR은 새 송금자에서 더 큰 차이로 뒤처져요. 엔도스랭크와 견주면, 두 기간 모두 새 허락 쌍에서 C-PR이 더 낮아요. 봄의 두 라벨 모두에서 C-PR의 계수는 허락 쪽 무게(가중치)가 커질수록 올라가요. 그래서 기간 밖에서는, 스펜더에 대해 송금 층이 허락 층이 주는 것 말고 더 보태는 것이 없어요.

빼 보기 실험(ablation, 한 부분을 바꾸어 무엇이 효과를 내는지 보는 실험)과 층의 크기가 두 부분을 모두 설명해요. S-PR(보증 씨앗 페이지랭크)은 송금 걷기를 그대로 두되, 허락 걷기의 결과에서 다시 출발해요. S-PR은 흐름 라벨(새 송금자 라벨)에서는 송금 걷기를 그대로 되풀이하고, 보증 신호는 거의 다 잃어요. 걷기가 어디로 갈지를 송금 화살표가 정할 때에는, 보증에서 얻은 순간 이동(걷다가 갑자기 다른 점으로 뛰는 일) 벡터가 바꾸는 것이 별로 없어요. C-PR에서는 허락 화살표가 걷기의 일부를 맡아요. 그래서 스펜더 홀드아웃에서 신호의 일부가 살아남아요. 하지만 송금 층이 훨씬 더 빽빽하고, $\lambda$는 두 층 모두에 나가는 화살표가 있는 소수의 점에서만 작용해요(\ref{sec:coupled-results}절). 섞은 점수는 송금 걷기를 이길 만큼은 허락 신호를 지켰지만, 허락 걷기를 따라잡기에는 너무 적게 지켰어요. C-PR의 층들은 금액을 그대로 담고, 모든 점에서 고르게 다시 출발해요. 반면 AWP는 송금 하나하나를 거의 한 번씩으로 세고, 활발한 송신자에서 다시 출발해요. 그리고 AWP는 두 기간 모두에서 새 송금자를 C-PR보다 더 잘 줄 세워요. 자료에서 배우거나 점마다 달라지는 $\lambda$, 엔도스랭크와 AWP처럼 무게를 단 층, 그리고 두 정상 벡터를 나중에 합치는 결합 방식은 시험하지 않았어요(\ref{ch:conclusion}장).

두 부분 모두 결과로 쳐요. 둘째 부분의 규칙은 결합 점수가 하나도 없을 때 정했어요. 그래서 $\lambda$를 그 규칙에 맞추어 조정했을 수가 없어요. 그리고 이것은 누구나 떠올릴 다음 단계에 한계를 그어요. 이동 행렬(전이 행렬) 단계에서 두 화살표 종류를 볼록 혼합(더하면 1이 되는 비율로 섞기)으로 섞으면, 어느 구성개념에서도 더 나은 한 층을 이기는 점수가 나오지 않아요. 첫째 부분의 규칙은 라벨을 꺼내기 전에 등록했고, 세 조건 가운데 하나에서 실패해요. 재현이 실제로 보여 주는 것은 더 좁아요. 송금 걷기에 허락 화살표를 더하면 권한 부여(대신 쓰도록 허락하는 일)에 대한 정보를 얻고, 이 정보는 새 자료에서도 다시 나타나요. 하지만 들어오는 흐름에서 잃는 것이 어느 정도를 넘지 않는다는 것은 보이지 못했어요. 그리고 두 기간 모두에서, 그냥 센 수가 두 스펜더 라벨을 어떤 걷기보다도 더 잘 줄 세워요.

\dissertationSubheading[sec:objectives-achievement]{연구 목표 달성}

\ref{ch:introduction}장은 다섯 가지 목표마다 재는 방법과 기준을 하나씩 붙여 두었어요. 표~\ref{tab:objectives-achievement}은 그 목표들을 증거와 맞추어 점검해요.

\begin{table}[htbp]
\centering
\caption{연구 목표, \ref{ch:introduction}장에서 밝힌 기준, 그리고 결과.}
\label{tab:objectives-achievement}
\fitwidth{%
\begin{tabular}{p{0.06\linewidth}p{0.36\linewidth}p{0.42\linewidth}p{0.10\linewidth}}
\toprule
목표 & 기준 & 증거 & 결과 \\
\midrule
O1 & 두 방법 사이 $\tau$(켄달의 타우, 두 줄 세우기가 얼마나 같은 순서인지 나타내는 수)의 95\% 구간에 $0$도 $1$도 들어 있지 않아요 & 표~\ref{tab:tie-sensitivity}의 마지막 줄: 구간 $[0.342,0.379]$가 끝에 닿지 않고 $(0,1)$ 안에 들어 있어요 & 달성 \\
O2 & 어떤 지표 가족에서 점수의 구간이 $0$을 벗어나면, 그 점수는 그 구성개념과 맞아요. 점수가 자기 화살표의 구성개념과 맞는 것은 의도한 구성개념 점검이에요 & 표~\ref{tab:alignment-hybrid}: 송금 그래프를 걷는 모든 점수가 송금 지표 가족에서($0.494$부터 $0.612$까지), 허락 그래프를 걷는 두 점수가 허락 지표 가족에서($0.375$와 $0.355$) 맞아요. EndorseRank의 계수는 동점이 정해요(표~\ref{tab:tie-sensitivity}) & 달성 \\
O3 & 실행 시간 비율을 화살표 수 비율과 함께, 하나로 정해 둔 측정 절차 아래에서 보고해요 & 표~\ref{tab:benchmark-runtime}, \ref{tab:benchmark-scaling}와~\ref{tab:benchmark-scaling-incohort}: 설정마다 평균, 표준편차(SD), 메모리, 반복 횟수, $|V|$, $|E|$ & 달성 \\
O4 & 구간이 $0$을 벗어나면 그 점수는 라벨을 예측해요. 페이지랭크는 짝지은 차이의 구간이 $0$보다 위에 있을 때만 자기 차수보다 나아요 & 표~\ref{tab:holdout-spenders}과~\ref{tab:holdout-diff}: 새 허락 쌍에서의 C-PR ($\lambda=0$)만 빼면 모든 점수가 두 라벨을 예측해요. 차수와의 차이는 $-0.233$(C-PR, $\lambda=1$), $-0.317$(EndorseRank), $-0.127$(AWP)이고, 모두 구간이 $0$ 아래예요. 6월부터 8월까지의 기간에서도 다시 $0$ 아래예요(표~\ref{tab:fresh-posthoc}) & 달성(첫째 부분); 달성 못 함(둘째 부분) \\
O5 & (a) 9월 14일의 규칙: 홀드아웃에서 C-PR은 새 허락 쌍에서 C-PR ($\lambda=1$)보다 나쁘지 않고, 새 송금자에서 C-PR ($\lambda=0$)보다 나아야 해요. 같은 기간 안에서는, 송금 지표 가족과 허락 지표 가족에서 C-PR의 점 추정값(값 하나로 어림한 값)이 앞서는 층의 구간 안에 있어야 하고, 시빌 안정성(가짜 주소를 많이 만드는 공격의 영향을 덜 받게 잰 안정성)에서는 두 층을 모두 넘어야 해요. (b) 9월 28일에 쓰고 10월 1일에 등록한 규칙: 6월부터 8월까지의 기간에서 C-PR은 새 허락 쌍에서 C-PR ($\lambda=0$)을 이기고, 새 송금자와 트레이더의 청산 없이 닫은 비율에서는 0.02보다 더 많이 나쁘지 않아야 해요(95\% 구간의 아래 끝이 $-0.02$보다 위) & (a) 표~\ref{tab:holdout-diff}: 홀드아웃의 두 조건이 모두 실패해요. 같은 기간: 송금($0.527$)은 C-PR ($\lambda=0$)의 구간 안에 있고, 허락($-0.001$)은 C-PR ($\lambda=1$)의 구간에서 한참 벗어나 있어요(표~\ref{tab:alignment-hybrid}). 시빌 안정성은 두 층을 모두 넘어요(표~\ref{tab:tau-diff-hybrid}). 허락 부분이 실패해요. (b) 표~\ref{tab:fresh-contrasts}: 새 허락 쌍 $+0.179$ $[0.132,0.225]$와 청산 없이 닫은 비율 $+0.018$ $[0.011,0.024]$는 조건을 채워요. 새 송금자 $-0.011$ $[-0.025,0.003]$는 아래 끝이 $-0.02$보다 낮아서 실패해요 & (a) 달성 못 함; (b) 달성 못 함 \\
\bottomrule
\end{tabular}
}
\end{table}

목표 밖의 발견 네 가지도 결론의 모양을 정해요. 이 발견들은 주장에 한계를 긋는 역할을 해요. 매칭 코호트에서 엔도스랭크의 같은 기간 계수는 동점에 따라 모양이 정해져요(\ref{sub:finding-rq2}절). 확장 지갑 풀 실험은 코호트 그래프보다 큰 그래프에서는 한 번도 돌지 않았어요(\ref{sub:finding-rq3}절). 매칭 코호트의 트레이더에서는 어떤 점수도 이익을 내며 닫은 비율을 줄 세우지 못해요(\ref{sub:trader-results}절). 그리고 봄 스펜더에서는 금액을 그대로 써서 걸은 허락 층이 두 라벨을 엔도스랭크보다 더 잘 줄 세웠어요. 그러니 이 구성개념에서는 허락에 무게를 어떻게 주는지가 중요해요(\ref{sub:threat-construct}절). 원시 차수가 모든 페이지랭크보다 더 잘 예측한다는 것은 이 네 가지에 들지 않아요. 그것은 O4의 둘째 부분이기 때문이에요.

\dissertationSubheading[sec:theoretical-implications]{이론에 주는 뜻}

이 결과는 이론의 세 영역과 관련이 있어요. 금융 그래프에서 퍼뜨림이 뜻하는 것, 온체인(블록체인 위의) 점수의 구성 타당도, 그리고 가명성(진짜 이름 대신 주소로만 활동함) 아래에서 믿음의 경제학(믿음에 드는 비용과 얻는 것)이에요.

\dissertationSubsubheading[sub:theory-propagation]{송금을 넘어선 퍼뜨림의 뜻}

페이지랭크는 처음에 하이퍼링크(웹 페이지끼리 잇는 링크)를 위해 만들어졌고, 이때 하이퍼링크를 인용 관계의 망으로 읽었어요\cend{15}. Do와 Do\cend{4}는 이 읽는 방식을 이더리움(Ethereum) 거래로 옮겨 왔어요. AWP\cend{3}는 여기에 금액 무게와 나이(오래된 정도) 무게, 그리고 활발한 송신자에서 다시 출발하기를 더해 다듬었어요. 엔도스랭크는 이 방식을 ERC-20(토큰 공통 규칙) 허락에 써요. 허락에서는 스펜더가 자기 대신 일하게 해 주는 주인이, 쓸 권한에 대해 방향이 있는 보증을 내주는 셈이에요\cend{18}. RQ1에서 본 순위 갈라짐은, 같은 점들 위에서도 화살표의 뜻을 바꾸면 순위가 바뀐다는 것을 보여 줘요. 그래서 그래프에 바탕을 둔 평판은 여러 방법이 모인 가족이고, 각 방법은 자기 화살표가 담은 관계에 비추어서만 타당해요. 개인끼리 직접 잇는 망(P2P)의 신뢰 시스템과 웹의 신뢰 시스템은 반대쪽에서 같은 결론에 닿았어요. EigenTrust\cend{16}와 TrustRank\cend{17}에서는 씨앗 집합과 화살표의 정의가 결과를 대부분 정해요. 그쪽의 화살표는 메시지를 주고받는 오버레이 망(인터넷 위에 따로 짠 연결망)이나 하이퍼링크 그래프에 속했지만, 엔도스랭크의 화살표는 토큰 표준에 원래 들어 있는 것이에요.

홀드아웃은 퍼뜨림이라는 비유가 가리는 단서 하나를 보태요. 화살표 종류를 바꾸자 점수가 미리 내다볼 수 있는 것도 바뀌었어요. 두 기간 모두에서, 새 허락에서는 허락 화살표로 만든 점수들이 앞섰고 송금만으로 만든 점수들은 낮게 머물렀어요. 화살표를 따라 퍼뜨리는 일은 여기에 보탠 것이 없어요. 그냥 센 허락 수는 어떤 걷기보다도 새 허락을 더 잘 예측했고, 그냥 센 송금 입차수는 어떤 걷기보다도 새 송금자를 더 잘 예측했어요. 이미 거래 상대가 많은 지갑이 계속 더 얻는 곳이라면 예상할 수 있는 일이에요. 이 자료에서 가치는 어떤 화살표를 고르느냐에 있고, 걷기는 대부분 개수를 매끄럽게 다듬을 뿐이에요. 결합은 한 층이 다른 층보다 훨씬 듬성듬성할 때 무슨 일이 생기는지 보여 줘요. 빽빽한 층이 길을 정하고, 스펜더 홀드아웃에서 듬성듬성한 층의 신호는 일부만 살아남아요.

\dissertationSubsubheading[sub:theory-construct]{여러 구성개념 영역과의 일치도}

점검들은 송금 중심성(송금 그물에서 얼마나 가운데에 있는지)과 허락으로 맡기기(위임)를 나누고, 모든 점수를 두 쪽 모두에 견주어 보도록 설계했어요. 각 점수는 자기가 걷는 화살표의 구성개념과 잘 맞고, 섞은 점수들은 더 빽빽한 층과 잘 맞아요. 고전적인 구성 타당도 이론은, 서로 다른 구성개념을 재는 도구라면 이런 모양이 나올 것이라고 예측할 거예요\cend{22}. Messick\cend{33}의 통합된 관점은 한 걸음 더 나아가서, 점수를 썼을 때 생기는 결과까지 그 점수의 타당도의 일부로 세요. 순위 상관(두 순위가 함께 움직이는 정도)\cend{20,21}은 분류 문턱값(어디서부터 나눌지 정하는 값) 없이도, 순서만 뜻이 있는 평판 점수에 이 논리를 실제로 적용하게 해 줘요. 그리고 짝지은 부트스트랩(다시 뽑기)\cend{34}은 이 모양을 시험할 수 있게 해 줘요. 같은 이론은 이런 모양이 결정할 수 없는 것도 말해 줘요. 두 도구가 서로 다른 구성개념을 잴 때, 어떤 점검에서 앞서는 쪽은 바로 그 점검을 위해 만든 쪽이에요. 그러니 계수로 도구들의 순서를 매겨도, 어느 쪽이 평판을 더 잘 재는지는 알 수 없어요. 엔도스랭크의 동점은 실제적인 점 하나를 보태요. 두 순위 상관은 얻을 수 있는 가장 큰 값이 같을 때만 견줄 수 있어요.

\dissertationSubsubheading[sub:theory-pseudonymity]{값싼 가명 아래의 평판}

누구나 볼 수 있는 장부에서는 신분을 싸게 만들 수 있어요\cend{5}. 그리고 이더리움의 계정 방식에서는 새 주소를 만들 때 등록 단계가 아예 필요 없어요\cend{36}. 그래서 들어오는 송금이나 허락에 상을 주는 평판 시스템은 모두 시빌 전략(가짜 주소를 많이 만들어 속이는 방법)에 노출돼요. 시빌 전략에서는 공격자가 자기가 움직이는 주소들끼리 화살표를 만들어요\cend{11}. 실제 자료에서 얻은 발견은 이것을 배제하지 못해요. 다만 관찰한 표본에서 본 활동에 허락 그래프와 송금 그래프가 서로 다르게 반응한다는 것은 보여 줘요. \ref{sec:sybil-cost-model}절에서 우리는 허락 그래프를 공격하는 비용을 수식으로 정리해요. 서로 허락해 주는 주소들로 이루어진 닫힌 농장(공격자가 만든 가짜 주소 무리)은 바깥 스펜더가 필요 없어요. 또 송금 자전 거래(자기끼리 사고팔아 거래가 많은 척하기)와 달리 토큰 잔고도 필요 없어요. 그래서 이 농장의 화살표에는 가스(거래 수수료) 말고는 돈이 거의 들지 않아요.

\dissertationSubheading[sec:practical-implications]{실제에 주는 뜻과 활용 시나리오}

서비스(프로토콜)와 기반 시설 제공자는 송금 허브를 찾아내는 점수와, 권한 부여를 살피는 평판 점수 가운데에서 골라야 해요. 잰 결과는 점수 하나를 추천하기보다, 상황(시나리오)에 따라 점수를 고르는 쪽을 뒷받침해요.

\dissertationSubsubheading[sub:deploy-transfer]{송금 허브와 흐름 감시}

어떤 서비스는 지금의 흐름을 그려 내는 일을 가장 먼저 챙겨요. 시장 조성자(늘 사고팔 값을 내놓아 거래를 돕는 쪽)를 알아보는 일, 거래 경로 분석, 송금에 바탕을 둔 규정 준수(법과 규칙을 지키는 일) 어림 규칙이 그런 목적이에요. 이런 서비스는 그 흐름을 들어오는 화살표 수(입차수)와 받은 양의 합에서 곧바로 읽을 수 있어요. AWP\cend{3}는 여기에 퍼뜨림, 시간 무게(오래된 일일수록 가볍게 세는 무게), 활발한 송신자에서 다시 출발하기를 더해요. AWP는 지금 들어오는 흐름에 가장 가까운 걷기이고(표~\ref{tab:alignment-hybrid}), 송금 그래프를 걷는 점수 가운데에서는 나중의 새 송금자에 가장 가까워요(표~\ref{tab:holdout-spenders}). 하지만 기간 밖에서는, 두 기간 모두에서 그냥 센 송금 입차수가 나중의 새 송금자를 AWP보다 더 잘 줄 세웠어요. 흐름을 미리 내다보고 싶은 서비스는 개수에서 시작해야 해요.

\dissertationSubsubheading[sub:deploy-allowance]{권한 부여를 살피는 빠른 갱신}

쓸 권한을 분명하게 주는 일이 가장 중요한 곳도 있어요. 토큰 금고(vault), 허락이 있어야 쓸 수 있는 연동 기능, 허락을 믿음의 행동으로 보는 위험 관리 엔진이 그래요. 이런 곳에서는 서비스가 허락 화살표를 쓸 까닭이 있어요. 가장 싼 형태는 스펜더의 받은 허락 수예요. 받은 허락 수는 어떤 페이지랭크보다도 새 허락을 더 잘 예측했고(표~\ref{tab:holdout-diff}), 풀이 프로그램(solver)이 아예 필요 없어요. 엔도스랭크는 이것을 퍼뜨린 형태예요. 허락 일치도(표~\ref{tab:alignment-hybrid}), 홀드아웃 결과(표~\ref{tab:holdout-spenders}), 짧은 풀이 시간(표~\ref{tab:benchmark-runtime}) 덕분에, 엔도스랭크는 보증해 주는 주인이 누구인지까지 반영한 순위를 원하는 운영자에게 괜찮은 선택이에요. 하지만 이 논문은 이 성질이 어떤 라벨이라도 더 낫게 한다는 것을 보이지 못했어요. 또 모든 점에서 고르게 순간 이동할 때, 퍼뜨린 점수는 닫힌 농장이 싸게 부풀릴 수 있어요(\ref{sec:sybil-cost-model}절). 운영자는 권한 부여의 뜻과 속도를 얻는 대신, 송금 정보를 거의 담지 않은 점수를 받아들이는 거예요.

\dissertationSubsubheading[sub:deploy-hybrid]{평가 설계와 함께 쓰기}

두 신호를 합칠지는 실제로 쓰는 일의 문제이기도 하고 방법의 문제이기도 해요. 그리고 증거는 이 둘에 서로 다르게 답해요. 운영자는 권한 부여를 살피는 꾸준한 감시에는 엔도스랭크를, 정해진 때마다 하는 송금 허브 점검에는 AWP를 돌릴 수 있어요. 대신 두 번 풀어야 하는 비용이 들어요. 두 화살표 종류를 하나의 걷기로 접어 넣는 것은 이 자료에서 득이 되지 않았어요. C-PR은 AWP보다 비용이 더 들고(표~\ref{tab:benchmark-runtime}), 두 기간 모두에서 새 허락에서는 엔도스랭크에, 새 송금자에서는 AWP에 뒤져요(\ref{sub:finding-rq5}절). 그러니 마음대로 고를 수 있는 운영자는 두 점수를 모두 계산하고 따로 두어야 해요. 금액을 그대로 무게로 쓴 송금 걷기 하나에 묶인 운영자라면, 허락 화살표를 더한 대가로 권한 부여 정보를 얻었고, 이 정보는 새 자료에서도 다시 나타났어요. 하지만 들어오는 흐름에서 $0.02$보다 큰 손실이 생길 가능성은 배제할 수 없었어요. 그리고 AWP를 쓰는 운영자는 여기서 만든 그대로의 C-PR로 바꾸면 새 송금자 순위에서 손해를 볼 거예요. 처리 과정이 공개되어 있으니, 허락 로그(계약이 남기는 기록)가 있는 ERC-20 체인이라면 어디서든 이 분석을 다시 할 수 있어요(부록~\ref{ch:appendix-eval}). 실시간 위험 오라클(바깥 정보를 블록체인에 알려 주는 장치)과 연결하는 일은 앞으로의 연구로 남겨 둬요(\ref{ch:conclusion}장).

\dissertationSubsubheading[sub:deploy-industry]{업계 도입 시나리오}

두 점수는 모두 순서를 알려 주는 신호이고, 다른 사람이 공개 자료에서 다시 계산할 수 있어요. 아래 시나리오는 각 점수가 누구에게 쓸모 있고 누구에게 쓸모없는지 말해 줘요.

\begin{itemize}
\item 탈중앙 거래소(DEX) 운영자. 지갑마다 수수료 등급이나 감시 순서를 정하는 서비스에는 자기 사용자에 대한 점수가 필요해요. 그런데 그 사용자들은 허락을 해 주는 주인이에요. 엔도스랭크는 그 사용자들 거의 모두에게 같은 값을 줘요(\ref{sec:rank-divergence}절). 그래서 이 목적에는 쓸 수 없어요. 송금에 바탕을 둔 점수나 단순한 활동 횟수는 쓸 수 있어요. 대신 엔도스랭크는 그 탈중앙 거래소가 연결해 쓰는 라우터(거래를 알맞은 곳으로 보내 주는 계약)와 금고에 순위를 매길 수 있어요.
\item 빌려주기(대출) 서비스. 대출 시장은 담보를 줄이지 않고도 가진 돈을 더 많이 쓰이게 하고 싶어 해요\cend{41}. 그리고 2020년의 시장 충격은 위험 설정값을 늦게 바꾸면 치르는 비용을 보여 줬어요\cend{9}. 빌리는 사람에 대한 신호라면 주인에게 순위를 매겨야 하는데, 엔도스랭크는 그렇게 하지 않아요. 등록한 청산 라벨(\ref{sub:fresh-holdout}절)은 송금 걷기나 결합 걷기가 주인에 대해 무엇이라도 말해 주는지 시험해요. 이 점수들은 채무 불이행(빌린 돈을 갚지 못하는 일)을 예측하는 도구가 아니에요.
\item 지갑과 보안 도구. 사용자와 연동하는 개발자는 어떤 스펜더 계약이 널리 허락을 받았는지 알고 싶어 해요. 몇 명의 주인이 그 스펜더를 허락했는지, 또는 그 스펜더의 엔도스랭크를 보여 주는 배지는 그 맡김을 분명히 보여 줘요\cend{18,19}. 이 배지는 그래프에서 점 하나의 위치를 알려 줄 뿐이에요. 보안 감사도 아니고, 나쁜 계약을 막아 준다는 약속도 아니에요.
\item 유동성 공급자(거래에 쓸 돈을 맡겨 두는 사람)와 시장 조성자. 이 참여자들은 재고와 가격 차이(스프레드)를 관리하려고 송금 허브를 찾고, 흐름이 어떻게 짜여 있는지 봐야 해요. 이 목적에는 송금 횟수와 AWP가 알맞은 도구예요\cend{3}. 그 이점은 상대를 가려내는 비용이 낮아지는 것이지, 손익이 좋아지는 것이 아니에요.
\item 에어드롭(공짜 토큰 나눠 주기)과 거버넌스(투표로 하는 운영) 프로그램. 시빌 무리는 토큰 나눠 주기와 거버넌스 자격을 늘 악용해요\cend{11,32}. 여기서 정의한 엔도스랭크는 거의 지켜 주지 못해요. 닫힌 농장은 자기 인원 비율(전체 주소 가운데 농장 주소가 차지하는 비율)의 몇 배를 모으고, 그것을 목표 하나에 몰아줄 수 있기 때문이에요(\ref{sec:sybil-cost-model}절). 프로그램은 서로 다른 내력을 가진 주인들에게서 여러 갈래로 들어온 보증을 요구할 수 있어요. 닫힌 무리 안에서만 도는 몫을 깎아서 셀 수도 있어요. 또는 순간 이동과 매달린 점(나가는 화살표가 없는 점)의 몫을 검증한 씨앗 집합에만 보낼 수도 있어요. 하지만 이 가운데 어느 것도, 정직한 주인에게서 보증을 사 오는 공격자는 막지 못해요.
\item 분석 업체와 규정 준수 업체. 거래소, 펀드, 규정 준수 팀에는 바깥 사람도 다시 계산할 수 있는 거래 상대 순위가 필요해요. 엔도스랭크와 AWP는 정한 기간과 설정값으로 공개하고 감사받을 수 있어요. 이 점에서 회사만 가진 비공개 기계 학습 모델과 달라요\cend{2}. 다만 시빌 공격(가짜 주소를 많이 만들어 속이는 공격)에 얼마나 노출되는지 밝히고, 이익을 예측하는 점수처럼 내세우지 않아야 해요.
\end{itemize}

이 시나리오 가운데 어느 것에서도 평판이 담보를 대신하지는 않아요.

\dissertationSubsubheading[sub:deploy-multiples]{함께 쓰는 기반 시설로서의 표준 점수}

여러 쪽이 함께 쓰는 점수는, 저마다 그 점수를 다시 계산하고 어떻게 조작될 수 있는지 알 수 있을 때만 쓸모가 있어요. 엔도스랭크와 AWP는 공개된 처리 과정으로 공개 로그에서 다시 계산할 수 있어요(부록~\ref{ch:appendix-eval}). 그리고 \ref{sec:sybil-cost-model}절과 \ref{sec:regulatory-ethics}절은 두 점수가 어떻게 조작될 수 있는지 기록해요. 둘 다 이익이 날지를 내다본 값은 아니에요.

\dissertationSubheading[sec:societal-impact]{사회와 실제에 미치는 영향}

디파이(DeFi, 탈중앙 금융) 시장에서는 거래 상대를 꼼꼼히 살피는 데 돈이 많이 들고, 그럴 수 있는 힘도 고르게 나뉘어 있지 않아요. 돈이 많은 참여자는 송금 기록 전체를 분석하는 비공개 위험 도구에 돈을 낼 수 있을지도 몰라요. 반면 개인 참여자는 들은 이야기에 기대거나, 살피는 일을 아예 건너뛰어요. 사용자가 이미 내준 권한으로 만들고 다시 계산하기에 싼 점수(표~\ref{tab:benchmark-runtime})는, 그 권한을 달라고 하는 계약을 살피는 문턱을 낮춰 줘요. 사용자가 전에 내준 허락은 토큰이 빠져나가는 길 가운데 하나예요(\ref{sec:approval-security}절). 그래서 스펜더를 가운데에 둔 보증 점수는, 사용자가 직접 마주하고 송금 양으로는 잴 수 없는 위험에 대한 정보를 전해 줘요.

서비스 개발자에게 이 결과는, 하나의 “평판” 구성개념 대신 구체적인 점수들을 쓰자고 말해요. 각 점수는 자기가 돕는 결정에 맞춘 것이에요. 이렇게 하면 흐름 점수로 권한 부여를 결정하는 위험이 줄어요. 두 점수 모두 공개 자료에서 다시 계산할 수 있으니, 사용자도 자기가 얼마나 노출되어 있는지 직접 점검할 수 있어요.

더 넓게 보면, 여기서 연구한 점수들은 공개된 이벤트와 버전을 매긴 설정값으로 만들고, 모든 표를 그 입력에서 다시 만들 수 있어요. Packin과 Lev-Aretz\cend{2}가 다룬 탈중앙 신용 점수 시스템은 속이 보이지 않아요. 그와 달리 이런 점수는, 디파이 서비스를 분석할 때 요구해야 한다고 규제 기관이 제안하는 설명 정보를 줘요\cend{61}. 이런 점수에는 의무도 따라와요. \ref{sec:regulatory-ethics}절에서 설명하는 것처럼, 그래프에 바탕을 둔 순서 통계를 신분이나 신용 기록이나 사람 증명(주소 뒤에 진짜 사람 한 명이 있다는 증명)으로 잘못 알지 않도록 해야 해요.

실험은 블록체인 하나에서 6개월 동안 했어요. 그래서 결과는 보증 구성개념에 대해 말해 줄 뿐, 신용 위험에 대해서는 말해 주지 않아요. 그리고 사회적 가치가 있다면, 그것은 결과를 그 구성개념 안에서 쓰는지에 달려 있어요. 가치를 어떻게 나눌지에 영향을 주는 평판 지표는 무엇이든 악용될 수 있어요. 다음 절에서는 엔도스랭크를 조작하는 데 얼마가 드는지 모형으로 나타내요.

\dissertationSubheading[sec:sybil-cost-model]{허락 그래프에서의 시빌 공격 비용 모형}

시빌 공격은 값싼 가명을 이용해 믿을 만하다는 거짓 인상을 만들어요\cend{11}. 송금 그래프에서는 나쁜 참여자가 자기가 움직이는 계정들끼리 자산을 옮겨서, 입차수와 들어오는 흐름을 인위적으로 부풀릴 수 있어요. 그 대가로 거래 수수료와, 어쩌면 재고 손실을 치러요\cend{29}. 허락 그래프는 다른 공격 면을 드러내요. 그곳의 화살표는 하나하나가 실제 송금이 아니라 권한 부여를 기록해요. 이 절에서 우리는 가짜 허락으로 엔도스랭크를 부풀리는 간단한 비용 모형을 세워요.

\dissertationSubsubheading[sub:sybil-setup]{공격 설정}

공격자가 주소 집합 $\mathcal{S}$를 움직인다고 해요. 그리고 보증 그래프 $G_{\text{approve}}=(V,E,w)$에 있는 정직한 주소들을 $\mathcal{H}=V\setminus\mathcal{S}$라고 써요. 무게가 $w_{uv}>0$인 화살표 $(u,v)$는 $u$가 $v$에게 자기 토큰을 써도 된다고 권한을 주었다는 기록이에요. 무게는 식~\eqref{eq:approve-weight}의 무게예요. 토큰마다 기간 안의 마지막 허락의 시간 무게에 그 금액의 상한이 있는 값 무게(금액이 아무리 커도 1을 넘지 않는 무게)를 곱하고, 이것을 모든 토큰에 대해 더해요. 감쇠 계수를 $d\in(0,1)$로 두면, 엔도스랭크는 $G_{\text{approve}}$ 위의 페이지랭크 벡터 $\mathbf{r}$을 다음 식의 해로 구해요\cend{15}.
\begin{equation}
\label{eq:pagerank-sybil}
\mathbf{r} = d\, P^{\top} \mathbf{r} + \frac{1-d}{|V|}\mathbf{1},
\end{equation}
여기서 $P$는 식~\eqref{eq:transition}의 이동 행렬이에요. 나가는 화살표가 있는 점은 자기 줄을 나가는 허락 무게에 비례하게 나눠요. 매달린 점, 곧 나가는 화살표가 없는 점은 자기 줄의 모든 칸이 똑같이 $1/|V|$ 값을 가져요. 그래서 그 점의 몫은 모든 점에 고르게 퍼져요\cend{23}. 이 고정점(되풀이해도 더는 바뀌지 않는 값)은 식~\eqref{eq:pagerank}의 반복이 닿는 고정점과 같아요. 구현에서는 바로 이 반복을 돌리고, 매달린 점의 몫을 고른 순간 이동 벡터를 거쳐 되돌려 보내요. 공격자의 목적은 목표 집합 $\mathcal{G}\subseteq\mathcal{S}$의 평균 엔도스랭크를 가장 크게 만드는 것이에요.
\begin{equation}
\label{eq:objective}
J(\mathcal{G}) = \frac{1}{|\mathcal{G}|}\sum_{i\in\mathcal{G}} r_i,
\end{equation}
다만 공격 비용에 대한 예산 제약을 지켜야 해요. 이 비용은 바로 다음에 정의해요.

\dissertationSubsubheading[sub:sybil-cost]{비용 구성 요소}

공격자는 화살표 집합 $E_{\mathcal{S}}$를 만들어요. 이 화살표들은 모두 $\mathcal{S}$에서 출발해요. 또 정직한 주인에게서 $\mathcal{S}$로 들어오는 화살표 집합 $E_{\mathcal{H}\to\mathcal{S}}$를 살 수도 있어요. 비용은 관찰 기간 $q$개에 걸쳐 세요. 이 수는 목표 주소들이 순위를 지키도록 하려는, 점수 새로 고침의 횟수예요. 모든 화살표에 드는 비용은 가스뿐이에요. 다른 구성 요소는 특정한 화살표나 전략에서만 생겨요.

\begin{enumerate}
\item 가스 비용. \texttt{approve}나 \texttt{permit}을 부를 때마다 가스가 들어요. 허락을 한 번 바꿀 때 드는 수수료의 기댓값을 $c_{\mathrm{gas}}$라고 써요. 엔도스랭크는 관찰 기간 안의 \texttt{Approval} 이벤트만 읽어요. 그래서 어떤 화살표가 한 기간의 그래프에 들어가려면, 그 허락이 그 기간 안에 기록되어야 해요. $E_{\mathcal{S}}$를 기간 $q$개 동안 그대로 두려면, 공격자는 기간마다 모든 허락을 다시 내보내야 해요. 그러면 시간 무게도 가장 큰 값인 $\sigma(0)=0.86$ 가까이에 머물러요. 그래서 가스를 $q$번 내야 해요.
\begin{equation*}
C_{\mathrm{gas}} = q\,|E_{\mathcal{S}}|\,c_{\mathrm{gas}}.
\end{equation*}

\item 노출 비용. 허락은 토큰을 묶어 두지 않아요. 그리고 주인은 허락할 때 그 토큰을 조금도 가지고 있지 않아도 돼요. 금액도 거의 상관없어요. 값 무게 $V$는 원래 단위(소수점 조정을 하지 않은 가장 작은 단위) 몇 개짜리 허락에도 무제한 허락과 거의 같은 무게를 줘요. 그리고 줄 정규화(한 줄의 합이 1이 되게 나누기)를 하고 나면, 나가는 화살표가 하나뿐일 때 그 화살표는 무게가 얼마든 $P_{uv}=1$이 돼요. 그래서 $\mathcal{S}$ 안의 화살표는 아무것도 노출하지 않고, 가스 말고는 비용이 없어요. 노출은 토큰을 가진 공격자 주소가 스펜더 $v\notin\mathcal{S}$를 허락할 때만 생겨요. 예를 들어 잘 알려진 계약을 쓰는 평범한 사용자처럼 보이려고 그럴 수 있어요. 그런 허락을 받으면 $v$가 그 토큰을 옮길 수 있기 때문이에요. 이런 화살표는 점수의 몫이 $\mathcal{S}$ 밖으로 빠져나가게도 해요. 기간 하나 동안 이런 노출 하나가 공격자에게 주는 비용의 기댓값을 $\ell_{uv}$라고 하면 다음과 같아요.
\begin{equation*}
C_{\mathrm{exp}} = q\sum_{(u,v)\in E_{\mathcal{S}},\ v\notin\mathcal{S}} \ell_{uv},
\end{equation*}
닫힌 농장에서는 이 값이 0이에요.

\item 허락 취소와 들킴 비용. 정상적인 사용자는 가끔 오래된 허락을 거두어요. 기간마다 허락이 바뀌지 않고 똑같이 다시 나타나는 농장은 눈에 띌 수도 있어요. 들키는 일의 비용과, 들키지 않으려고 공격자가 하는 허락 취소의 비용을 합친 기간당 기댓값을 $c_{\mathrm{rev}}$라고 해요. 기간 $q$개 동안에는 다음과 같아요.
\begin{equation*}
C_{\mathrm{rev}} = q\,c_{\mathrm{rev}}.
\end{equation*}

\item 사 온 정직한 화살표. 공격자는 정직한 주인에게 돈을 주고 자기 주소를 허락하게 할 수도 있어요. 예를 들어 사용자에게 권한을 달라고 하는 앱을 통해서 그럴 수 있어요. $h\in\mathcal{H}$이고 $a\in\mathcal{S}$인 사 온 화살표 $(h,a)$에도 같은 기간 규칙이 적용되고, 기간마다 다음 비용이 들어요.
\begin{equation}
\label{eq:marginal}
c_{\mathrm{in}}(h,a) = c_{\mathrm{gas}} + p_h + \rho_h,
\end{equation}
여기서 $p_h$는 정직한 주인에게 주는 돈이에요. $\rho_h$는 공격자가 그 허락으로 토큰을 써 버릴 위험에 대비해 주인이 더 달라고 하는 위험 웃돈(위험 프리미엄)이에요.
\end{enumerate}

기간 $q$개 동안 공격의 전체 비용은 다음과 같아요.
\begin{equation}
\label{eq:total-cost}
C = C_{\mathrm{gas}} + C_{\mathrm{exp}} + C_{\mathrm{rev}} + q\sum_{(h,a)\in E_{\mathcal{H}\to\mathcal{S}}} c_{\mathrm{in}}(h,a).
\end{equation}
닫힌 농장에서는 $C_{\mathrm{exp}}=0$이고 $E_{\mathcal{H}\to\mathcal{S}}$가 비어 있어요. 그래서 $C=q\,(|E_{\mathcal{S}}|\,c_{\mathrm{gas}}+c_{\mathrm{rev}})$가 돼요.

\dissertationSubsubheading[sub:sybil-bound]{영향력과 비용 대비 효과의 한계}

페이지랭크의 민감도에 대해 흔히 알려진 결과에 따르면, 화살표 하나가 $\mathbf{r}$을 움직이는 정도는 그 화살표가 출발하는 점의 정상 분포 몫과, 정규화한 나가는 무게에 비례해요\cend{23}. 먼저 닫힌 농장을 생각해 봐요. $\mathcal{S}$의 모든 점에는 나가는 화살표가 적어도 하나 있고, 이 화살표들은 모두 $\mathcal{S}$ 안에서 끝나요. 그리고 어떤 정직한 주인도 $\mathcal{S}$ 안의 주소를 허락하지 않아요. 매달린 점들의 집합을 $D$라고 해요. 이때 매달린 점은 모두 $\mathcal{H}$에 있어요. 그리고 $r_D=\sum_{j\in D} r_j$를 매달린 점들의 정상 분포 몫이라고 해요. 식~\eqref{eq:pagerank-sybil}에 따르면, 점 $i\in\mathcal{S}$는 농장에서 $d\sum_{j\in\mathcal{S}} r_j P_{ji}$를, 순간 이동에서 $(1-d)/|V|$를, 매달린 점의 줄에서 $d\,r_D/|V|$를 받아요. 농장 점들의 줄은 $\mathcal{S}$ 안에서 합이 1이에요. 그래서 $\mathcal{S}$ 전체에 대해 더하면 $\sum_{i\in\mathcal{S}} r_i = d\sum_{i\in\mathcal{S}} r_i + |\mathcal{S}|\,(1-d+d\,r_D)/|V|$를 얻어요. 곧 다음과 같아요.
\begin{equation}
\label{eq:closed-cluster}
\sum_{i\in\mathcal{S}} r_i = \frac{|\mathcal{S}|}{|V|}\cdot\frac{1-d+d\,r_D}{1-d}.
\end{equation}
첫째 인수(곱해진 두 부분 가운데 앞의 것)는 농장의 인원 비율이고, 둘째 인수는 농장이 그 비율의 몇 배를 모으는지 말해 줘요. 둘째 인수는 $r_D=0$일 때만, 곧 매달린 점이 없는 그래프에서만 1이에요. 그렇지 않으면 매달린 점의 줄이 정직한 쪽 몫의 일부를 농장에 넘겨줘요. 매달린 점들은 $\mathcal{S}$ 밖에 있으니 $r_D\le 1-\sum_{i\in\mathcal{S}} r_i$예요. 이것을 식~\eqref{eq:closed-cluster}에 넣으면, 어떤 그래프에서나 성립하는 한계를 얻어요.
\begin{equation}
\label{eq:cluster-mass}
\sum_{i\in\mathcal{S}} r_i \le \frac{|\mathcal{S}|}{(1-d)\,|V| + d\,|\mathcal{S}|}.
\end{equation}
$|V|$에 비해 작은 농장이라면, 이 한계는 인원 비율의 $1/(1-d)$배에 가까워요. 매칭 코호트의 엔도스랭크 그래프에서, 처리 과정에 들어 있는 페이지랭크 코드로 계산하면 $d=0.75$, $0.85$, $0.95$일 때 $r_D=0.532$, $0.553$, $0.572$예요. 그래서 이 그래프의 닫힌 농장은 자기 인원 비율의 2.6배, 4.1배, 11.9배를 모아요. 같은 코드로 직접 점검해 봐도 결과가 맞아요. 저마다 다음 주소를 허락하는 주소 20개짜리 고리를 그래프에 더하면, $d=0.85$에서 이 고리는 자기 인원 비율의 4.09배를 모아요.

식~\eqref{eq:cluster-mass}은 목표 하나의 점수가 아니라 농장 전체의 몫을 제한해요. 그리고 공격자는 그 몫을 한곳에 몰 수 있어요. 별 모양을 생각해 봐요. 농장 주소 $m$개가 저마다 목표 $t$를 허락하고, $t$는 그 주소들을 하나하나 다시 허락해요. 모든 농장 주소 $a$는 나가는 이웃이 $t$ 하나뿐이에요. 그래서 $P_{at}=1$이고, $t$의 줄은 농장 위에서 합이 1이에요. 각 점이 순간 이동과 매달린 점의 줄에서 받는 몫을 $b=(1-d+d\,r_D)/|V|$라고 써요. 그러면 균형 식 $r_a=b+d\,r_t P_{ta}$와 $r_t=b+d\sum_a r_a$에서 다음을 얻어요.
\begin{equation*}
r_t = \frac{b\,(1+dm)}{1-d^2}.
\end{equation*}
$d=0.85$에서 $b=9.62\times10^{-5}$이고, 이 값은 허락을 하나도 받지 않은 모든 지갑이 똑같이 받는 공통 점수예요. $m=10$일 때(주소 11개, 그리고 기간마다 허락 20개), 별을 더한 그래프에서 목표는 $r_t=3.27\times10^{-3}$에 이르러요. 이것은 코호트 지갑 가운데 가장 높은 엔도스랭크($4.23\times10^{-4}$)의 7.7배예요. 코호트 지갑은 대부분 주인이에요. 점수는 $m$에 따라 일정한 크기씩(직선으로) 늘어요. 농장 주소를 하나 더할 때마다 $r_t$에 $d\,b/(1-d^2)$가 더해지고, 기간마다 허락이 2개씩 더 들어요.

그러니 코호트의 모든 지갑보다 높이 올라가는 데 정직한 주인에게서 들어오는 화살표는 필요 없어요. 하지만 사 온 화살표는 얻는 것을 더 키워요. 농장에서 정직한 주소로 나가는 화살표가 없을 때, 정직한 주인 $h$에게서 사 온 화살표 $(h,a)$는 1차 근사(가장 큰 부분만 남겨 어림한 값)로 다음만큼을 농장의 몫에 더해요.
\begin{equation*}
\Delta\sum_{i\in\mathcal{S}} r_i \approx \frac{d\,r_h\,P_{ha}}{1-d}
\end{equation*}
여기서 $P_{ha}$는 새 화살표가 $h$의 줄에서 가져가는 비율이에요. 얻는 것은 허락한 금액이 아니라 이 비율에 달려 있어요. $V$는 원래 단위 몇 개 이상인 허락을 모두 거의 한 번씩으로 세요. 그래서 $h$에게서 산 새 허락은 얼마를 허락하든 $h$의 줄에서 $\sigma(0)/(o(h)+\sigma(0))$만큼을 가져가요. 이것은 대략 ($h$의 허락 수 + 1)분의 1이에요. 그래프에서 중앙값인 주인에게 이 비율은 $0.30$이에요. 그래서 공격자가 살 화살표는 다른 허락이 적은 주인의 허락이고, 되도록 $r_h$가 높은 주인의 것이 좋아요. 그리고 그 허락은 별 의미 없을 만큼 적은 토큰 금액이어도 돼요. 이런 허락은 거의 아무것도 노출하지 않아요. 그래서 식~\eqref{eq:marginal}의 위험 웃돈 $\rho_h$는 0에 가까울 수 있어요. C-PR의, 금액을 그대로 쓴 허락 층에서는 반대예요. 그곳에서는 무제한 허락만이 주인의 줄에서 큰 비율을 가져가요. 그리고 무제한 허락은 그 토큰의 주인 잔고 전부를 공격자 마음대로 쓸 수 있게 해요(\ref{sub:threat-construct}절). 공격의 비용 대비 효과는 다음과 같아요.
\begin{equation}
\label{eq:efficiency}
\eta = \frac{\Delta J(\mathcal{G})}{C},
\end{equation}
여기서 $\Delta J(\mathcal{G})$는 공격 화살표가 만들어 내는, 식~\eqref{eq:objective}의 목적값이 늘어난 양이에요. 이 양은 식~\eqref{eq:total-cost}의 $C$가 값을 치르는 기간 $q$개 동안 유지돼요.

식~\eqref{eq:closed-cluster}--\eqref{eq:efficiency}은 여기서 정의한 엔도스랭크를 닫힌 농장이 싸게 부풀릴 수 있다는 것을 보여 줘요. 순간 이동이 고르고 매달린 점의 몫도 고르게 다시 나눌 때, 닫힌 무리는 자기 인원 비율에 묶여 있지 않아요. 그리고 위의 두 결과 모두 정직한 주인에게서 들어오는 화살표가 하나도 필요 없어요. 닫힌 무리를 정말로 묶어 두는 장치는 걷기 자체를 바꿔요. TrustRank\cend{17}처럼 순간 이동과 매달린 점의 몫을 검증한 씨앗 집합에만 보내면, 그 집합 밖의 무리는 정직한 주인에게서 오는 화살표를 통해서만 몫을 얻어요. 공격자는 그런 화살표를 식~\eqref{eq:marginal}의 값을 치르고 사야 해요. 송금 자전 거래에는 수수료와 재고가 들지만, 닫힌 허락 농장은 토큰을 갖고 있지 않고 가스 말고는 거의 돈을 내지 않아요. AWP의 다시 출발하기 규칙도 송금 쪽을 지켜 주지 못해요. 농장 주소들이 서로 적은 금액을 보내면 그 주소들의 활발함(활동 정도)이 올라가고, 그래서 다시 출발하는 몫을 자기 쪽으로 끌어와요. 시빌 보정 안정성 대리 지표(\ref{sub:proxy-sybil}절)는 몇 안 되는 주소에서 되풀이되는 송금은 잡아내지만, 새로 만든 주소들로 된 농장은 잡아내지 못해요. 그래서 이 지표들은 공격을 가정한 평가를 대신할 수 없어요. 그런 평가는 앞으로의 연구로 남겨 둬요.

\dissertationSubheading[sec:regulatory-ethics]{규제와 윤리에서 생각할 점}

온체인 평판 지표는 공정함의 문제뿐 아니라 투명함의 문제도 일으켜요. 앞선 법 연구는, 정산이 온체인에서 이루어지더라도 탈중앙 금융 시장이 여전히 공시(정보를 밝히는 일) 규칙과 소비자 보호 규칙에 묶여 있다는 것을 보여 줬어요\cend{61}. Packin과 Lev-Aretz\cend{2}는 그들 나름의 경고를 보태요. 모델을 해석하거나 따질 수 없을 때, 또는 정당한 절차 없이 신분 정보 데이터베이스와 이어질 때, 탈중앙 신용 점수는 속을 볼 수 없는 상자(블랙박스)가 주는 해를 계속 이어 갈 수 있다는 거예요. 엔도스랭크와 AWP는 둘 다 공개된(오픈 소스) 알고리즘으로 온체인 이벤트에서 결정적으로(늘 같은 결과가 나오게) 계산해요. 그래서 누구나 오픈 소스 소프트웨어로 다시 계산할 수 있어요(부록~\ref{ch:appendix-eval}). 이 공개성은 블랙박스라는 비판에 어느 정도 답해요. 하지만 분배에 대한 걱정은 그대로 남겨요. 거래 활동이 적은 주소는, 오프체인(블록체인 밖의) 신용도가 어떻든 점수가 더 낮게 나오기 때문이에요.

윤리적으로, 이 점수들을 신분을 대신하는 지표로, 전통적인 뜻의 신용도 척도로, 또는 한 주소가 한 사람에게 속한다는 증거로 써서는 절대 안 돼요\cend{11}. 이 점수들은 정해진 그래프 구조 안에서 한 주소가 어디에 있는지를 순위로 알려 주는 지표예요. 어떤 서비스가 점수 문턱값을 근거로 금융 서비스를 내준다면, 점수 뒤에 있는 그래프 정의, 기간, 그리고 한계를 분명히 기록해야 해요. 한계 가운데에서도 특히 시빌 공격에 대한 노출을 기록해야 해요. 해석할 수 있음을 중심에 둔 규제 틀이라면, Kotb\cend{14}가 검토한 것 같은 비공개 기계 학습 모델보다 이런 순위를 더 좋아할 거예요. 다만 제공자가 대출 상환율을 얼마나 잘 예측하는지 부풀려 주장하지 않아야 해요\cend{2}.

열린 블록체인에서 사생활 보호는 장부의 성질 그 자체 때문에 한계가 있어요\cend{5}. 그리고 평판 점수를 계산할 때는 공개 자료 말고는 아무것도 쓰지 않아요. 그래서 윤리에서 가장 중요한 질문은 결과를 어떻게 쓰느냐예요. 각 주소 뒤에 있는 사람이 누구인지 더 알아내려는 시도는 권하지 않아요. 그리고 점수를 지문처럼 영원히 따라다니는 식별표로 다루어서는 안 돼요. 이 연구는 지갑 주소만 다루고, 그 주소를 실제 세상의 사람과 이으려고 하지 않아요(\ref{sec:ethics}절).

\dissertationSubheading[sec:threats-validity]{타당성을 위협하는 것}

여기서 보고한 발견은 아래의 한계에 비추어 읽어야 해요. 한계는 구성 타당도, 내적 타당도(연구 안의 절차 때문에 결론이 흔들리지 않는지), 외적 타당도(바깥 자료로도 맞는지), 공격에 대한 타당도(누가 일부러 속이려 해도 결과가 버티는지)로 묶었어요.

\dissertationSubsubheading[sub:threat-construct]{구성 타당도}

\begin{itemize}
\item 이 점수는 신용 등급이 아니에요. 담보 없는 대출에서 나온 채무 불이행 라벨은 분석에 하나도 들어가지 않아요.
\item 같은 기간 점검은 지갑 바로 곁에서 센 허락 통계와 송금 통계에 기대요. 홀드아웃에서 라벨은 앞으로 생길 새 허락과 앞으로 생길 새 송금자예요.
\item 홀드아웃의 구성개념. 앞으로 생길 새 송금자가 여전히 흐름 사건이듯, 앞으로 생길 새 허락도 여전히 보증 사건이에요. 둘 다 점수 기간 밖에 있지만, 어느 것도 따로 떨어진 신용 구성개념이나 거래 구성개념은 아니에요. 또 둘 다 새로 도착한 것의 개수라서, 지금의 거래 상대 수를 따라가는 점수라면 무엇이든 유리하게 만들어요(\ref{sub:finding-rq4}절).
\item 분석 단위. 엔도스랭크는 허락을 받는 쪽에 점수를 매겨요. 매칭 코호트는 대부분 주인으로 이루어져 있고, 엔도스랭크는 주인들을 동점으로 남겨 둬요. 그래서 그 코호트에서 엔도스랭크의 같은 기간 계수는 몇 안 되는 지갑과, 허락 그래프 안과 밖으로 지갑을 가르는 일에 기대요. 그리고 엔도스랭크의 기간 밖 증거는 스펜더에서 나오는데, 그 스펜더는 대부분 계약이에요(\ref{sub:cohorts}절). 트레이더의 평판에 대한 주장은 탐색적으로 돌린 트레이더 실험과 등록한 기간에 기대요.
\item 자기 구성개념과의 일치는 어느 정도 저절로 생겨요. 엔도스랭크는 스펜더의 받은 허락 수를 그래프 전체에 걸쳐 매끄럽게 다듬고, AWP는 들어오는 송금 수에 대해 똑같이 해요. 그래서 각 점수는 만드는 방식 때문에 자기 지표 가족과 맞을 것으로 예상돼요. 표~\ref{tab:tau-diff}의 점수 안 대비(한 점수의 두 계수를 견준 것)는 이 때문에 계산했어요. 하지만 엔도스랭크에서는 각 대비에 든 두 계수의 천장이 서로 달라요(\ref{sec:same-window-alignment}절). 그래서 그 대비에서 차이가 없다는 결과가 나와도 알려 주는 것이 없어요. 또 이 대비들은 그래프 밖의 지갑에 점수를 어떻게 주는지에 따라서도 달라져요(표~\ref{tab:tie-sensitivity}).
\item 값 무게. 엔도스랭크와 AWP는 각 금액을, 원래 단위에서 $b=1$로 둔 $V(z)$에 넣어 바꿔요. 원래 단위 10개 이상인 금액이면 이 값은 언제나 1과의 차이가 $10^{-4}$ 안이에요. 그래서 무게는 한 화살표가 허락이나 송금을 몇 번 담는지와 그것이 얼마나 최근인지를 기록하고, 금액에 대해서는 거의 아무것도 기록하지 않아요. 소수 자릿수(소수점 아래 자리 수)가 18인 토큰을 원래 단위 몇 개만큼 허락해도 무제한 허락처럼 세요. 그 논문(AWP를 발표한 논문)은 $b$의 값을 알려 주지 않아요. 그리고 가격으로 맞춘 금액에 맞추어 정한 값, 곧 크기가 중요해지게 하는 값은 시험하지 않았어요. C-PR의 금액을 그대로 쓴 층들은 반대쪽 끝에 있어요. 허락 그래프의 주인 대부분이 내준 무제한 허락이 주인의 줄을 거의 다 차지해요. 그리고 소수 자릿수가 다른 토큰의 금액을 단위 그대로 견주어요(\ref{sub:coupled-graph}절). 봄 스펜더 홀드아웃에서는 금액을 그대로 쓴 허락 걷기가 두 라벨을 엔도스랭크보다 더 잘 줄 세워요(표~\ref{tab:holdout-spenders}). 그러니 이 구성개념에서는 금액에 무게를 어떻게 주는지가 중요하고, 어느 무게 주기 방식도 옳은 방식이라고 보이지 않았어요. 이런 무게 선택은 $\mathbf{r}$과 순위 상관에까지 퍼져 들어가요\cend{23}.
\end{itemize}

\dissertationSubsubheading[sub:threat-internal]{내적 타당도}

\begin{itemize}
\item 함께 쓰는 관찰 기간. 모든 방법과 대리 지표는 2025년 12월 1일부터 2026년 5월 31일까지의 자료로 계산해요. 이 기간에만 있었던 시장 상황(가격 출렁임, 보상 프로그램, 서비스 변경)이 일치도를 올리거나 내릴 수 있어요. 감쇠 계수와 상위 토큰 부분 그래프에 대한 강건성 점검(조건을 바꿔도 결과가 버티는지 보기)은 설정값(하이퍼파라미터)과 토큰 쏠림에 대한 민감성을 낮춰 줘요. 하지만 그 가운데 어느 것도 달력 기간을 바꾸지는 않아요.
\item 코호트 고르기. 매칭 코호트는 아비트럼 원에서 기간 안에 포지션을 적어도 세 번 닫은 GMX V2 지갑 5,521개예요. 이 목록은 자료를 꺼내기 전에 GMX 자료로 정해 두었어요. 결과는 이 거래 집단 밖의 지갑에는 그대로 들어맞지 않을 수 있어요. 또 그래프가 코호트 지갑을 모두 보는 것도 아니에요(\ref{sub:cohorts}절). 그래프 밖의 지갑을 외톨이 점으로 남겨 두면, 엔도스랭크의 같은 기간 계수는 크게 바뀌지만 AWP의 계수는 전혀 바뀌지 않아요(표~\ref{tab:tie-sensitivity}). 그리고 표본 정의에 대한 민감도 분석은 평가 대상 집단이 바뀔 때 계수가 얼마나 크게 반응하는지 보여 줘요(표~\ref{tab:robustness-sample-size}).
\item 기준선 충실도(기준선이 원래 방법을 얼마나 잘 따랐는지). AWP는 화살표 무게와 다시 출발하기에서 발표된 방법을 따라요(\ref{sub:awp-graph}절). 그 논문은 설정값을 알려 주지 않아요. 그래서 $k$, $t_0$, $b$는 이 연구가 고른 값이에요. 그리고 원래 단위에서 $b=1$이면 값 무게는 계단(어떤 값을 넘으면 한꺼번에 뛰어오르는 모양)에 가까워요(위를 보세요). 그 논문은 이더리움 거래에 점수를 매겼고, 여기서는 AWP가 아비트럼 원의 ERC-20 송금에 점수를 매겨요. 엔도스랭크는 AWP의 무게를 바꾸지 않고 가져오고, 모든 점에서 고르게 다시 출발해요. 이 선택은 이 자료에서 두 가지 다시 출발하기 규칙을 견주어 본 뒤에 했어요(표~\ref{tab:endorserank-restarts}).
\item 분석 순서. 엔도스랭크와 AWP는 같은 기간 점검과 봄 라벨을 안 뒤에 점수를 매겼어요. 엔도스랭크는 6월부터 8월까지의 라벨도 안 뒤에 매겼어요. 그래서 두 점수가 들어간 비교는 모두 구간과 함께 보고해요. 그리고 C-PR을 AWP와 견주는 등록한 민감도 분석 말고는, 어느 비교도 미리 정해 두지 않았어요. C-PR을 그 송금 층 하나와 견주는 비교는 봄 라벨을 본 뒤에 골랐어요. 그래서 봄 기간에서는 이 비교를 탐색적인 것으로 다루고, 6월부터 8월까지의 기간에서 결정해요. 이 기간은 로그를 꺼내기 전인 2026년 10월 1일에 등록했어요(\ref{sub:fresh-holdout}절).
\item 시간을 잰 환경. 모든 실행 시간은 컴퓨터 한 대에서 잰 실제 걸린 시간이에요(부록~\ref{ch:appendix-eval}). 표~\ref{tab:benchmark-runtime}의 표준편차는 한 가지 설정을 되풀이해 푼 것만 다뤄요. 그리고 감쇠 계수를 바꿔 가며 돌린 실험에서는 그래프--감쇠 계수 쌍마다 시간을 한 번만 쟀어요. 컴퓨터 사이의 차이는 재지 않았어요. 그래서 비교는 초가 아니라 비율로 나타내요.
\end{itemize}

\dissertationSubsubheading[sub:threat-external]{외적 타당도}

\begin{itemize}
\item 체인 하나와 기간 하나. 증거는 아비트럼 원 하나에서, 6개월 동안 나왔어요. 이 발견을 다른 롤업(여러 거래를 묶어 처리하는 블록체인), 레이어 1 체인(기본 블록체인), 또는 대출 중심의 거래 장소로 넓히려면 여러 서비스에 걸친 검증이 필요해요. 디파이 거래 장소들은 서비스 설계\cend{8}와 시장 구조\cend{1}가 서로 다르기 때문이에요.
\item 풀 규모 키우기 대 일치도 규모 키우기. 확장 지갑 풀이 뒷받침하는 것은 매칭 코호트 그래프의 부분 그래프에서의 실행 시간에 대한 말뿐이에요(\ref{sub:runtime-scaling}절). 같은 기간 일치도는 여전히 매칭 코호트에 묶여 있고, 홀드아웃은 스펜더에 묶여 있어요. 효율이 좋아졌다고 해서, 집단 크기가 아무리 달라져도 같은 $\tau$ 모양이 유지된다고 받아들이면 안 돼요.
\end{itemize}

\dissertationSubsubheading[sub:threat-adversarial]{공격에 대한 타당도}

\begin{itemize}
\item 시빌과 자전 거래. 공격을 노린 허락 패턴이나 송금 패턴으로 순위가 조작될 수 있어요\cend{11}. 시빌 보정 대리 지표는 이 위험을 일부만 다뤄요. 그리고 \ref{sec:sybil-cost-model}절의 비용 모형은 수식으로 따진 논증이지, 실제로 공격해 보는 레드팀(일부러 공격해 보는 팀) 평가가 아니에요.
\item 전략적인 허락 행동. 엔도스랭크를 실제로 쓰기 시작하면, 이익을 따지는 참여자들은 순위를 모으려고 전략적으로 허락을 내줄 수 있어요. 실제로 쓰기 전에 관찰 자료로 어림한 $\tau$는 순위가 그런 행동을 얼마나 잘 버티는지 부풀려 보여 줄 수 있어요.
\end{itemize}

\dissertationSubheading[sec:limitations]{한계}

모든 결과는 아비트럼 원에서 6개월이라는 한 기간 동안 얻었어요. 대상은 매칭 지갑 코호트와 스펜더 홀드아웃 코호트였고, 홀드아웃에서는 매칭 코호트의 트레이더 일부 묶음도 썼어요. 엔도스랭크의 강점은 이래요. 허락 점검과 맞고, AWP와 다른 순위를 주고, 그래프가 듬성듬성한 덕분에(희소성) 비용이 낮아요. 그리고 스펜더에서는 두 기간 모두 새 허락 쌍을 줄 세우는데, 이 라벨에서 송금만으로 만든 점수들은 낮게 머물러요. 약점은 이래요. 순위에 송금 정보가 거의 없고, 같은 기간 계수가 자기 동점에 따라 모양이 정해지고, 홀드아웃 계수가 자기가 다듬는 원시 받은 허락 수보다 낮아요. 그리고 허락을 받지 않는 트레이더에 대해서는 할 말이 거의 없어요.

결합 연산자(두 층을 섞어 걷는 계산 규칙)는 $\lambda$의 정해진 세 값에서 시험했어요. 이때 순간 이동은 고르게 했고, 이동 행렬 단계에서 점마다 볼록 혼합으로 섞었고, 금액을 그대로 무게로 쓴 두 층을 썼어요. 엔도스랭크와 AWP처럼 무게를 단 층, 점마다 다르거나 자료에서 배운 무게, 이동 행렬의 줄 대신 정상 벡터를 섞는 방식, 그리고 다른 마르코프 연쇄(다음 걸음이 지금 있는 점에만 달린 무작위 걷기)는 시험하지 않았어요. 그래서 RQ5에 대한 두 답은 이 자료에서 그 방법 가족에만 해당해요. 값 설정값 $b$의 다른 값도 시험하지 않았어요(\ref{sub:threat-construct}절). 두 번째 단계의 자료 꺼내기가 끝나지 않았기 때문에, 확장 지갑 풀 실험은 매칭 코호트 그래프보다 큰 그래프에서는 한 번도 돌지 않아요. 더 오래 걸리는 기준선 실행들, 곧 다른 페이지랭크 변형과 대출 활동을 중심으로 만든 그래프는 따로 보관해 두었어요(부록~\ref{ch:appendix-archive}). 이 자료에는 담보 없는 대출의 채무 불이행 라벨이 없어요. 그리고 실제로 쓰는 일은 여전히 \ref{sec:regulatory-ethics}절에서 생각한 점들에 묶여 있어요.

엔도스랭크는 오프체인 증거를 하나도 쓰지 않아요. 휴먼 패스포트(Human Passport) 같은 증명 도구는 신분 확인, 생체 정보, 다른 사람의 보증을 써서 한 주소가 서로 다른 한 사람에게 속한다는 것을 확인해요\cend{64}. 이런 도구와 달리, 엔도스랭크는 순위가 높은 두 주소를 서로 다른 사람이 움직인다는 것을 보장하지 않아요. 그래서 엔도스랭크는 이미 받아들여진 집단 안의 순서로 읽어야 하고, 시빌 방어 수단으로 읽으면 안 돼요. 엔도스랭크를 계산하기 전에 증명 관문으로 주소를 걸러 내는 것은 그럴듯한 쓰임새예요. 하지만 자격 증명 파밍(자격 증명을 가짜로 모으기)\cend{66} 같은 그 관문 자체의 약점이 그대로 넘어올 거예요. 그리고 이 조합을 평가하려면 시빌 라벨이 붙은 자료가 필요한데, 이 연구는 그런 자료를 모으지 않았어요.

이 한계들은 결론으로 말할 수 있는 것을 좁혀요. 이 코호트와 이 기간에서, 엔도스랭크는 스펜더를 위한, 싸고 허락에 특화된 평판 점수예요. 엔도스랭크는 신용 점수 방법이 아니고, 시빌 공격을 버티는 신분 확인 장치도 아니고, 서비스 차원의 담보를 대신하는 것도 아니에요\cend{9}. 엔도스랭크가 기대는 허락 화살표는 온체인 평판에 보탤 만한 가치가 있고, 받은 허락 수는 그 화살표를 쓰는 가장 싼 방법이에요. \ref{ch:conclusion}장에서는 무엇을 보였는지, 그 앎이 얼마나 가치 있는지, 그리고 이 한계를 넓히려면 무엇이 필요한지를 다시 정리해요.
